\chapter{先看懂问题：我们想预测什么}
\label{ch:overview}

\section{从一个水下充电场景说起}
设想两只线圈隔着海水传输交流电能。刚接通电源时，线圈铜线、绝缘层、封装和海水都接近环境温度。通电一段时间后，铜线因为电阻产生热量，海水内部也因感应电流而发热。热量从铜经过绝缘层传入封装，随后向海水扩散。

升温又影响电磁过程：铜的电导率随温度改变；海水温度影响海水电导率和感应电流；感应电流的磁场改变线圈的完整复阻抗。阻抗改变电流，电流再改变后续发热量。相同电压下，经历过不同通断电历史的系统可能具有不同的温度、阻抗和效率。

我们希望输入\textbf{线圈形状、相对位置、封装尺寸、介质物性、电路和电压历史}，随时输出：
\begin{itemize}
\item 两端口复阻抗 $Z_{11}(t),Z_{12}(t),Z_{22}(t)$、实际电流和传输效率；
\item 发射铜、接收铜、两只封装及海水的温度演化和热点；
\item 海水三维复电流密度 $\mathbf J(\mathbf x,t)$、焦耳热密度 $q(\mathbf x,t)$；
\item 启动、停止、冷却、再次启动以及输入改变后的连续响应。
\end{itemize}

\begin{plainbox}
\textbf{三个关键词。}
\textbf{物理直接法}在每一热步重新求解完整电磁和热问题，作为参考计算；
\textbf{降阶模型（ROM）}用少数模式表示复杂空间温度，继续求解传热方程；
\textbf{神经代理}根据几何和过去的热源快速预测动态响应，再与物理模型共同完成计算。
\end{plainbox}

\section{为什么需要两种时间尺度}
以 200\,kHz 为例，交流电的一个周期约 $5\,\mu$s，而明显升温往往需要数秒到数百秒。因此，系统不沿着每个电磁周期计算温度，而是采用\textbf{快电磁、慢热学}的两时间尺度：
\begin{enumerate}
\item 给定当前温度，在交流频率下求稳态相量电磁场、端口复阻抗和周期平均损耗；
\item 用热时间步推进温度，重新计算电磁阻抗和实际电流，再推进下一步。
\end{enumerate}
这种双向关系正是我们研究的核心：\textbf{电磁损耗导致升温，温度变化又改变后续电磁损耗}。

\section{系统全貌}
\begin{figure}[htbp]
\centering\resizebox{\textwidth}{!}{\begin{tikzpicture}[node distance=9mm and 10mm]
\node[drive,text width=12.2cm] (inp) {用户输入：线圈形状、相对位置、两只封装、海水与热参数、补偿电路、随时间变化的电压};
\node[phys,below=of inp,text width=12.2cm] (bf) {原 BFZI V150：冷态双端口阻抗、单位端口三维电流场、空间损耗、热算子};
\node[nerve,below left=12mm and -5mm of bf,text width=5.4cm] (net) {40 维几何条件 + 3 路功率历史\\分数阶 SOE 神经预测};
\node[phys,below right=12mm and -5mm of bf,text width=5.4cm] (tan) {8 个海水温度模态\\原 BFZI 物理切线校正};
\node[drive,below=30mm of bf,text width=12.2cm] (circuit) {合成总复阻抗 $\mathbf Z(t)$ $\longrightarrow$ 双端口电路 $\longrightarrow$ 实际电流 $\mathbf I(t)$};
\node[therm,below=of circuit,text width=12.2cm] (heat) {同源三维海水热源与铜损耗 $\longrightarrow$ 32 阶守恒热 ROM $\longrightarrow$ 三维温度 $T(\mathbf x,t)$};
\node[phys,below=of heat,text width=12.2cm] (out) {逐时刻输出：复阻抗、功率和效率、铜/双封装/海水温度、三维 $\mathbf J$ 与 $q$};
\draw[arr] (inp)--(bf);
\draw[arr] (bf.south) -- ++(0,-.14) -| (net.north);
\draw[arr] (bf.south) -- ++(0,-.14) -| (tan.north);
\draw[arr] (net.south) |- (circuit.west);
\draw[arr] (tan.south) |- (circuit.east);
\draw[arr] (circuit)--(heat);
\draw[arr] (heat)--(out);
\draw[fb] (heat.west) -- ++(-1.0,0) |- (net.west);
\draw[fb] (heat.east) -- ++(1.0,0) |- (tan.east);
\node[font=\footnotesize,text=BFgreen,rotate=90] at ([xshift=-12mm]heat.west) {下一热步反馈};
\end{tikzpicture}
}
\caption{BFZI 电磁热神经代理的主要模块及温度反馈。}
\label{fig:system}
\end{figure}
请从图~\ref{fig:system} 的上方向下读。首先确定几何和材料；原 BFZI 提供冷态物理信息；神经记忆网络预测动态海水阻抗变化，物理灵敏度对它作高精度校正；电路确定两只线圈的实际电流；电流决定空间热源；热 ROM 更新空间温度，再把温度送回下一时间步。

系统区分两个工作阶段。\textbf{离线阶段}为某个几何构建冷态阻抗、单位端口三维场、热算子、热 ROM 以及必要的阻抗温度导数，训练神经网络并保存模型。\textbf{在线阶段}只需加载这些文件，输入新电压或负载历史，即可快速进行多步演化，避免每步执行完整 BFZI Green/GMRES 电磁求解。

\section{各模块在做什么}
\begin{table}[htbp]\centering\small
\caption{模块分工概览}
\begin{tabularx}{\textwidth}{p{3.3cm}X}
\toprule
模块 & 用通俗语言说明 \\
\midrule
BFZI 电磁基础 & 告诉我们线圈如何在海水中产生场、端口阻抗和发热。\\
冷态物理数据 & 为新几何保存初始温度下的物理基准。\\
分数阶 SOE 网络 & 把过去的功率输入压缩成短期与长期记忆，预测电磁温变趋势。\\
32 阶热 ROM & 用 32 个温度状态代替数千个热节点，保留真实的材料传热。\\
海水温度灵敏度 & 用少量原 BFZI 电磁导数修正动态海水复阻抗。\\
双端口电路 & 将完整阻抗与电压、负载组合，求两端口复电流和功率。\\
三维场重构 & 把温度基和单位端口场恢复到原空间位置，查询 $T/J/q$。\\
跨几何共享 & 借用已有几何的热模型和电磁温度导数，减少新几何部署时间。\\
\bottomrule
\end{tabularx}
\end{table}

\section{区分两种已经实现的神经路径}
当前仓库包含两种具体实现。其一是 \code{GeomFrac}：以 40 维几何条件、三路功率输入和记忆状态生成 14 维动态预测，和三维热 ROM、八模态物理校正共同工作。第~\ref{ch:neural} 章专门讲解它。

其二是 C 阶段的 \code{FracNet}：先由物理热 ROM 和电磁温度切线得到复阻抗预测，再用因果 SOE 网络学习剩余的六个实数误差。它们服务于相同的动态系统，网络输入、训练目标、在线路径有所不同。第~\ref{ch:training} 章会并列说明。

\section{如何理解目前的研发状态}
B 阶段热 ROM、C 阶段分数阶残差代理、D 阶段跨几何共享，已有各自的历史真实计算和独立留出验证。随后物理直接法在更强激励、更广热参数和新几何下发现材料域与近源积分精度问题，于是重新进行了冻结的 12 组 S5 五阶积分验收。

本文编写时，S5 的 12 组空间检查已经通过，其中 8 组完成空间与时间全部四组合；最终冻结 12/12 验收仍在进行。上述历史代理精度应与最新 S5 教师再作统一对照。实时进度请查阅上层 \code{PROJECT\_PROGRESS.md}。

初学者建议依次阅读电磁模型（第~\ref{ch:em} 章）、热 ROM（第~\ref{ch:heat} 章）、神经网络（第~\ref{ch:neural} 章）和训练方法（第~\ref{ch:training} 章），再学习物理校正与完整耦合。
